% part: intuitionistic-logic
% chapter: propositions-as-types
% section: type-preservation

\documentclass[../../../include/open-logic-section]{subfiles}

\begin{document}

\olfileid{int}{pty}{tpr}

\olsection{টাইপ সংরক্ষণ}

আলোচিত হ্রাস-রূপান্তর ও বিন্যাস-রূপান্তর
\Log{tN2i} এবং \Log{tN3i} ব্যবস্থার জন্যও
সরাসরি সংজ্ঞায়িত করা যায়। \Log{tN3i}-তে
প্রসঙ্গ~$\Gamma$-র সাপেক্ষে $\lambd$-পদ~$N$-এর
টাইপ হলো সেই !!{formula}~$!A$, যার জন্য
$\Gamma \Sequent N : !A$-এর !!a{proof} আছে।
$\lambd$-পদের $\beta$-হ্রাসবিধিগুলি !!{proof}s-এর
হ্রাস-রূপান্তর ও বিন্যাস-রূপান্তরের হুবহু
অনুরূপ হওয়ার একটি গুরুত্বপূর্ণ ফল আছে:
প্রসঙ্গ~$\Gamma$-র সাপেক্ষে $\lambd$-পদের
টাইপ $\beta$-রূপান্তরের পরেও একই থাকে।
এই ধর্মকে \emph{টাইপ সংরক্ষণ} বলে।

\begin{prop}
$\Gamma$-র সাপেক্ষে $N$-এর টাইপ~$!A$
এবং $N \redone N'$ হলে, $\Gamma$-র
সাপেক্ষে $N'$-এর টাইপও~$!A$।
\end{prop}

\begin{proof}
  $N$-এর গঠন এবং $\Gamma \Sequent N : !A$-এর
  !!{proof}s~$\delta$-র উচ্চতার উপর আরোহে
  নিম্নের কথাটি প্রতিষ্ঠিত হয়: $M$ যদি
  $N$-এর উপপদ হয়, তাহলে $\delta$-তে
  $\Gamma' \Sequent M : !B$-তে শেষ হওয়া
  একটি উপ-!!{proof}~$\delta_1$ আছে। $M$
  রিডেক্স হলে $\delta_1$-এ একটি রূপান্তর
  প্রযোজ্য। $\delta_1$-এ হ্রাস প্রয়োগে
  $\Gamma' \Sequent M' : !B$-এর
  !!a{proof}~$\delta_1'$ পাই।
  \Log{tN3i}-এর রূপান্তরগুলি দেখলে বোঝা যায়,
  $M'$ হলো $M$-এর সংকুচিত রূপ। $\delta$-তে
  $\delta_1$-এর স্থানে $\delta_1'$ বসালে
  $\Gamma \Sequent N' : !A$-এর
  !!a{proof}~$\delta'$ পাওয়া যায়; এখানে
  $N$-এর উপপদ~$M$-এর স্থানে~$M'$ বসিয়ে
  $N'$ তৈরি হয়েছে।

  যেমন, \Intro{\land}-এর পর \Elim{\land}
  প্রয়োগের হ্রাস-রূপান্তর দেখি:
  \[
  \bottomAlignProof
  \AxiomC{}
  \Deduce$\Gamma \fCenter N_1 : !B_1$
  \AxiomC{}
  \Deduce$\Gamma \fCenter N_2 : !B_2$
  \RightLabel{\Intro{\land}}
  \BinaryInf$\Gamma \fCenter \pair{N_1}{N_2} : !B_1 \land !B_2$
  \RightLabel{\Elim{\land}[i]}
  \UnaryInf$\Gamma \fCenter \proj{i}{\pair{N_1}{N_2}} : !B_i$
  \DisplayProof
  \quad\redone\quad
  \bottomAlignProof
  \AxiomC{}
  \Deduce$\Gamma \fCenter N_i : !B_i$
  \DisplayProof
  \]
  দেখা যায়, ডান দিকের !!{proof} অর্থাৎ
  $\delta_1'$-এর প্রমাণপদ $N_i$। বাঁ দিকের
  !!{proof}~($\delta_1$)-এর প্রমাণপদ
  $\proj{i}{\pair{N_1}{N_2}}$-তে
  $\beta$-রূপান্তর প্রয়োগের ফল ঠিক এটিই।

  অথবা \Intro{\lif}-এর পর \Elim{\lif}
  অনুমান এবং তার হ্রাস-রূপান্তর দেখি:
  \begin{multline*}
  \bottomAlignProof
  \AxiomC{$\delta_2$}
  \Deduce$x: !B_1, \Gamma \fCenter N : !B_2$
  \RightLabel{\Intro{\lif}}
  \UnaryInf$\Gamma \fCenter \lambd[\typeof{x}{!B_1}][N] : !B_1 \lif !B_2$
  \AxiomC{}
  \RightLabel{$\delta_3$}
  \Deduce$\Gamma \fCenter M : !B_1$
  \RightLabel{\Elim{\lif}}
  \BinaryInf$\Gamma \fCenter (\lambd[\typeof{x}{!B_1}][N] M) : !B_2$
  \DisplayProof
  \quad\redone\\
  \bottomAlignProof
  \AxiomC{}
  \RightLabel{$\delta_3$}
  \Deduce$\Gamma \fCenter M : !B_1$
  \RightLabel{$\Subst{\delta_2}{\delta_3}{x:!B_1}$}
  \Deduce$\Gamma \fCenter \Subst{N}{M}{x} : !B_2$
  \DisplayProof
  \end{multline*}
  এখানেও ডান দিকের !!{proof}-এর প্রমাণপদ
  বাঁ দিকের !!{proof}-এর প্রমাণপদ
  $\beta$-হ্রাস করলে যা হয়, ঠিক সেটিই।
  অবশ্য $\delta_2$-র উচ্চতার উপর আরোহে
  সতর্কভাবে যাচাই করতে হবে: $\delta_2$-তে
  $\Gamma' \Sequent x:!B_1$ রূপের সব
  স্বতঃসিদ্ধের স্থানে $\Gamma \Sequent
  M:!B_1$-এর !!{proof}~$\delta_3$ বসালে
  সত্যিই $\Gamma \Sequent
  \Subst{N}{M}{x} : !B_2$-এর
  !!a{proof}~$\Subst{\delta_2}{\delta_3}{x:!B_1}$
  পাওয়া যায়।
\end{proof}

!!{proof}s ও $\lambd$-পদের ঘনিষ্ঠ
সঙ্গতির আরেকটি ফল আছে। প্রসঙ্গ~$\Gamma$-র
সাপেক্ষে $N$ যদি টাইপ~$!A$-এর একটি
$\lambd$-পদ হয় এবং তাতে রিডেক্স~$M$
থাকে (অর্থাৎ সেটি স্বাভাবিক রূপে নেই),
তাহলে $\Gamma \Sequent N : !A$-এর
সংশ্লিষ্ট \Log{tN3i} প্রমাণ~$\delta$-তে
$\Gamma' \Sequent M: !B$-তে শেষ হওয়া
এমন উপ-!!{proof} আছে যাতে রূপান্তর
প্রয়োগ করা যায়। $M$-এর স্থানে তার
সংকুচিত রূপ বসিয়ে $N \redone N'$ হলে,
$\delta$-তে সংশ্লিষ্ট হ্রাস প্রয়োগের
ফলে $\Gamma \Sequent N' : !A$-এর
!!a{proof} পাওয়া যায়। অতএব স্বাভাবিক
রূপে নেই এমন প্রত্যেক $\lambd$-পদ অন্য
একটি পদে $\beta$-হ্রাস পায়। এই ধর্মের
নাম \emph{অগ্রগতি}।

অগ্রগতি ও টাইপ সংরক্ষণ মিলে নিশ্চিত করে
যে $\lambd$-পদের $\beta$-হ্রাস কখনো
“আটকে যায়” না: কোনো পদ সঠিকভাবে
টাইপযুক্ত এবং স্বাভাবিক রূপে না থাকলে,
সেটি একই টাইপের আরেকটি সঠিকভাবে
টাইপযুক্ত পদে হ্রাস পায়।

\end{document}
